\chapter{Atom Berelektron Banyak dan Simbol Suku}\label{ch:b3:many-electron-atoms}

Pada tahun 1868 sebuah garis kuning dalam spektrum Matahari mengungkap
unsur yang belum dikenal di Bumi, yaitu helium. Ketika helium akhirnya
diisolasi di laboratorium pada tahun 1895, spektrumnya tampak seperti
spektrum \emph{dua} unsur: dua keluarga garis, yang untuk sementara waktu
dianggap berasal dari ``ortohelium'' dan ``parahelium'', dan yang tidak
pernah saling bertukar cahaya. Helium hanya ada satu. Kedua keluarga itu
adalah \omterm{def:b3:many-electron-atoms:singlet-triplet}{keadaan triplet} dan \omterm{def:b3:many-electron-atoms:singlet-triplet}{keadaan singlet} dari atom yang sama, yang
dipisahkan oleh spin elektron dan oleh sebuah aturan yang lebih mendasar
daripada gaya apa pun: fungsi gelombang dua elektron harus berganti tanda
ketika kedua elektron itu dipertukarkan. Bab ini menelusuri aturan itu dari
prinsip Pauli sampai \omterm{def:b3:many-electron-atoms:term}{simbol suku} yang dipakai para spektroskopis untuk
menandai setiap keadaan setiap atom dan ion --- termasuk ion-ion logam
transisi yang warnanya dijelaskan dalam \cref{ch:b3:complex-spectra-magnetism}.

\begin{recall}
Jilid Tahun 1 menandai elektron dengan empat bilangan kuantum $n$, $l$,
$m_l$, $m_s$, menyusun konfigurasi dasar dengan prinsip Pauli, aturan
Klechkowski, dan aturan Hund, serta menghitung elektron tak berpasangan.
Jilid Tahun 2 memberikan energi orbital dan aturan Slater untuk penapisan.
\Cref{ch:b3:quantum-model-systems} menyediakan \omterm{def:b3:quantum-model-systems:operator}{operator}, \omterm{def:b3:quantum-model-systems:operator}{nilai eigen}, momentum
sudut rotor ($\hat L^2$ dengan \omterm{def:b3:quantum-model-systems:operator}{nilai eigen} $l(l+1)\hbar^2$), dan atom hidrogen.
\end{recall}

\section{Spin dan spin-orbital}

Elektron membawa momentum sudut intrinsik, yaitu spinnya, dengan bilangan
kuantum $s = \frac12$: $\hat S^2$ mempunyai satu \omterm{def:b3:quantum-model-systems:operator}{nilai eigen} saja, $s(s+1)\hbar^2 =
\frac34\hbar^2$, dan $\hat S_z$ dua \omterm{def:b3:quantum-model-systems:operator}{nilai eigen} $m_s\hbar = \pm\frac12\hbar$.
Kedua fungsi spin ditulis $\alpha$ ($m_s = +\frac12$) dan $\beta$
($m_s = -\frac12$); keduanya ortonormal, $\langle\alpha|\alpha\rangle =
\langle\beta|\beta\rangle = 1$ dan $\langle\alpha|\beta\rangle = 0$. Spin tidak
mempunyai padanan klasik; spin bukanlah rotasi sebuah bola kecil, dan spin
sama sekali tidak muncul dalam Hamiltonian nonrelativistik.

\begin{definition}[Spin-orbital elektron]\label{def:b3:many-electron-atoms:spin-orbital}
Yang disebut \emph{spin-orbital}\index{spin-orbital} adalah hasil kali
sebuah orbital ruang $\phi(\mathbf r)$ dan sebuah fungsi spin: $\phi\alpha$
atau $\phi\beta$. Sebuah spin-orbital memuat deskripsi lengkap satu elektron.
\end{definition}

Dengan demikian, orbital atom yang ditandai $(n, l, m_l)$ dalam jilid Tahun 1
memberikan dua \omterm{def:b3:many-electron-atoms:spin-orbital}{spin-orbital}; itulah sebabnya sebuah subkulit menampung
$2(2l+1)$ elektron. Pertanyaan yang sebenarnya adalah mengapa sebuah
\omterm{def:b3:many-electron-atoms:spin-orbital}{spin-orbital} menampung paling banyak satu elektron.

\section{Elektron yang tak terbedakan dan determinan Slater}

\begin{definition}[Partikel tak terbedakan, fungsi gelombang antisimetris]\label{def:b3:many-electron-atoms:indistinguishable}
Partikel-partikel disebut \emph{tak terbedakan}\index{partikel tak terbedakan}
jika tidak ada pengukuran yang dapat membedakannya: setiap elektron identik
dengan setiap elektron lain. Fungsi gelombang beberapa elektron disebut
\emph{fungsi gelombang antisimetris}\index{fungsi gelombang antisimetris} jika
fungsi itu berganti tanda ketika koordinat (ruang dan spin) dua elektron
mana pun dipertukarkan:
$\Psi(\dots,i,\dots,j,\dots) = -\Psi(\dots,j,\dots,i,\dots)$.
\end{definition}

Ketakterbedaan saja mensyaratkan $|\Psi|^2$ tidak berubah oleh pertukaran,
sehingga $\Psi$ paling banyak hanya dapat berganti tanda. Alam telah memilih:
fungsi gelombang setiap sistem elektron --- setiap sistem fermion --- bersifat
antisimetris. Ini adalah sebuah postulat mekanika kuantum, yang dikukuhkan
oleh setiap spektrum atom.

\begin{definition}[Determinan Slater]\label{def:b3:many-electron-atoms:slater}
Untuk $N$ elektron dalam $N$ \omterm{def:b3:many-electron-atoms:spin-orbital}{spin-orbital} yang berbeda, $\chi_1, \dots, \chi_N$,
yang disebut \emph{determinan Slater}\index{determinan Slater} adalah
\[
\Psi = \frac{1}{\sqrt{N!}}\begin{vmatrix}
\chi_1(1) & \chi_2(1) & \cdots & \chi_N(1)\\
\chi_1(2) & \chi_2(2) & \cdots & \chi_N(2)\\
\vdots & & & \vdots\\
\chi_1(N) & \chi_2(N) & \cdots & \chi_N(N)
\end{vmatrix},
\]
dengan baris $i$ memuat elektron $i$ dalam setiap \omterm{def:b3:many-electron-atoms:spin-orbital}{spin-orbital}. Determinan
itu disingkat $|\chi_1\chi_2\dots\chi_N|$.
\end{definition}

\begin{theorem}[Prinsip Pauli]\label{thm:b3:many-electron-atoms:pauli}
\omterm{def:b3:many-electron-atoms:slater}{Determinan Slater} bersifat antisimetris, dan determinan itu lenyap jika dua
elektron menempati \omterm{def:b3:many-electron-atoms:spin-orbital}{spin-orbital} yang sama. Oleh karena itu, tidak ada dua
elektron dalam sebuah atom yang dapat memiliki keempat bilangan kuantum yang
sama.
\end{theorem}

\begin{proof}
Mempertukarkan elektron $i$ dan $j$ berarti mempertukarkan dua baris
determinan, yang mengubah tandanya: antisimetri berlaku untuk setiap pasangan.
Jika dua \omterm{def:b3:many-electron-atoms:spin-orbital}{spin-orbital} sama, dua kolomnya sama dan determinannya nol: keadaan
semacam itu tidak ada. Dua elektron dalam atom yang sama dengan $n$, $l$,
$m_l$ dan $m_s$ yang sama akan menempati \omterm{def:b3:many-electron-atoms:spin-orbital}{spin-orbital} yang sama.
\end{proof}

\begin{example}[Keadaan dasar helium]\label{ex:b3:many-electron-atoms:helium-ground}
Dengan kedua elektron dalam $1s$,
\[
\Psi = \frac{1}{\sqrt2}\begin{vmatrix}1s\alpha(1) & 1s\beta(1)\\ 1s\alpha(2) &
1s\beta(2)\end{vmatrix} = 1s(1)\,1s(2)\cdot\frac{1}{\sqrt2}\big[\alpha(1)\beta(2)
- \beta(1)\alpha(2)\big].
\]
Bagian ruangnya simetris, bagian spinnya antisimetris: kedua spin
berpasangan, dengan spin total nol.
\end{example}

\section{Helium: tolakan Coulomb dan pertukaran}

Eksitasikan satu elektron helium ke $2s$. Empat determinan dapat disusun
dari $1s\alpha$, $1s\beta$, $2s\alpha$, $2s\beta$ dengan satu elektron di setiap
orbital. Setelah disusun ulang, keempatnya menjadi hasil kali satu fungsi
ruang dan satu fungsi spin:
\[
\Psi_\pm = \frac{1}{\sqrt2}\big[1s(1)2s(2) \pm 2s(1)1s(2)\big]\times(\text{fungsi
spin}).
\]
Fungsi ruang antisimetris $\Psi_-$ harus dipasangkan dengan fungsi spin
simetris, yang jumlahnya tiga: $\alpha(1)\alpha(2)$,
$\beta(1)\beta(2)$ dan $\frac{1}{\sqrt2}[\alpha(1)\beta(2) + \beta(1)\alpha(2)]$,
yaitu ketiga komponen $M_S = 1, 0, -1$ dari spin total $S = 1$. Fungsi
simetris $\Psi_+$ harus dipasangkan dengan satu-satunya fungsi spin
antisimetris, $S = 0$.

\begin{definition}[Keadaan singlet dan triplet]\label{def:b3:many-electron-atoms:singlet-triplet}
Keadaan dengan spin total $S = 0$ disebut \emph{keadaan singlet}\index{keadaan singlet};
keadaan dengan spin total $S = 1$, yang ketiga komponennya $M_S = 1, 0, -1$
mempunyai energi yang sama bila tidak ada medan magnet, disebut
\emph{keadaan triplet}\index{keadaan triplet}.
\end{definition}

\begin{definition}[Integral pertukaran]\label{def:b3:many-electron-atoms:exchange}
Untuk dua orbital $a$ dan $b$ serta tolakan elektron $e^2/4\pi\varepsilon_0
r_{12}$, yang disebut \emph{integral pertukaran}\index{integral pertukaran} adalah
\[
K_{ab} = \iint a(1)b(2)\,\frac{e^2}{4\pi\varepsilon_0r_{12}}\,b(1)a(2)\,\dd\tau_1\dd\tau_2,
\]
yaitu integral yang sama dengan tolakan elektron klasik $J_{ab}$ (dengan
$a(1)b(2)$ di kedua sisi), kecuali bahwa elektron-elektronnya ditukar pada
satu sisi. $K_{ab}$ bernilai positif dan tidak mempunyai padanan klasik.
\end{definition}

\begin{proposition}[Energi singlet dan triplet]\label{prop:b3:many-electron-atoms:helium}
Sampai orde pertama dalam tolakan elektron, keadaan-keadaan $1s2s$ helium
mempunyai energi $E_\pm = E_{1s} + E_{2s} + J \pm K$, dengan singlet ($+$)
berada di atas triplet ($-$) sejauh $2K$.
\end{proposition}

\begin{proof}
Energi tak terganggu kedua $\Psi_\pm$ adalah $E_{1s} + E_{2s}$ (orbital mirip
hidrogen dengan muatan inti 2). Koreksi orde pertama adalah \omterm{def:b3:quantum-model-systems:hermitian}{nilai harapan}
tolakan $\hat V_{12}$ dalam keadaan ternormalisasi:
$\frac12\langle 1s(1)2s(2) \pm 2s(1)1s(2)|\hat V_{12}|1s(1)2s(2) \pm
2s(1)1s(2)\rangle$. Kedua suku langsung memberikan $\frac12(J + J)$; kedua
suku silang memberikan $\pm\frac12(K + K)$, karena $\hat V_{12}$ simetris
terhadap elektron-elektronnya. Jadi $\Delta E = J \pm K$, dan fungsi-fungsi
spin, yang ternormalisasi dan tidak disentuh oleh $\hat V_{12}$, hanya
mengalikan dengan 1.
\end{proof}

Triplet berada lebih rendah meskipun Hamiltoniannya sama sekali tidak memuat
spin: fungsi ruangnya lenyap ketika $\mathbf r_1 = \mathbf r_2$, sehingga kedua
elektron saling menghindar dan tolak-menolak lebih lemah. ``Lubang Fermi'' ini
adalah isi fisis aturan Hund yang pertama.

\begin{example}[Integral pertukaran helium, terukur]\label{ex:b3:many-electron-atoms:K}
% ledger: lev:He.2s3S1, lev:He.2s1S0
Tingkat-tingkat $1s2s$ helium terletak pada \qty{159856}{cm^{-1}} (${}^3$S) dan
\qty{166277}{cm^{-1}} (${}^1$S) di atas keadaan dasar: singlet itu
\qty{6421}{cm^{-1}}, yaitu \qty{0.796}{eV}, lebih tinggi, dan $K(1s,2s) = \qty{0.398}{eV}$. Gambar di bawah memperlihatkan kedua keluarga itu.
\end{example}

\section{Kopling Russell--Saunders dan simbol suku}

Dalam atom ringan, momentum sudut orbital elektron-elektronnya dijumlahkan
menjadi momentum total $\mathbf L$, spin-spinnya menjadi spin total
$\mathbf S$, dan baru sesudah itu $\mathbf L$ dan $\mathbf S$ berkopling
lemah satu sama lain.

\begin{definition}[Kopling Russell--Saunders]\label{def:b3:many-electron-atoms:russell-saunders}
Dalam \emph{kopling Russell--Saunders}\index{kopling Russell--Saunders},
keadaan-keadaan sebuah konfigurasi ditandai dengan \emph{bilangan kuantum
momentum sudut orbital total}\index{bilangan kuantum momentum sudut orbital total}
$L$ ($\hat{\mathbf L}^2 = L(L+1)\hbar^2$, dengan $M_L = \sum m_l$),
\emph{bilangan kuantum spin total}\index{bilangan kuantum spin total} $S$
($M_S = \sum m_s$), dan \emph{bilangan kuantum momentum sudut
total}\index{bilangan kuantum momentum sudut total} $J$, yang mengambil
nilai-nilai $|L - S|, \dots, L + S$.
\end{definition}

\begin{definition}[Suku spektroskopi, multiplisitas spin, simbol suku]\label{def:b3:many-electron-atoms:term}
Yang disebut \emph{suku spektroskopi}\index{suku spektroskopi} adalah
himpunan $(2L+1)(2S+1)$ keadaan sebuah konfigurasi dengan $L$ dan $S$
tertentu. Bilangan $2S + 1$ disebut \emph{multiplisitas spin}\index{multiplisitas spin}
suku itu. Yang disebut \emph{simbol suku}\index{simbol suku} adalah
${}^{2S+1}L$, dengan huruf S, P, D, F, G, H untuk $L = 0, 1, 2, 3, 4, 5$;
keadaan dengan $J$ tertentu ditulis ${}^{2S+1}L_J$ --- misalnya
\termsym{3}{P}{2}.
\end{definition}

\begin{proposition}[Kulit tertutup]\label{prop:b3:many-electron-atoms:closed-shell}
Subkulit yang terisi penuh mempunyai $L = 0$ dan $S = 0$; subkulit itu tidak
menyumbang apa pun pada \omterm{def:b3:many-electron-atoms:term}{simbol suku}.
\end{proposition}

\begin{proof}
Dalam subkulit penuh, setiap $m_l$ muncul dua kali dan setiap $m_s = \pm\frac12$
muncul $2l+1$ kali, sehingga $M_L = 0$ dan $M_S = 0$; dan hanya ada satu cara
untuk mengisinya (satu determinan). Satu keadaan tunggal dengan
$M_L = M_S = 0$ hanya dapat termasuk dalam $L = 0$, $S = 0$.
\end{proof}

\begin{proposition}[Menghitung keadaan-keadaan sebuah konfigurasi]\label{prop:b3:many-electron-atoms:count}
$N$ elektron dalam subkulit dengan $2(2l+1)$ \omterm{def:b3:many-electron-atoms:spin-orbital}{spin-orbital} memberikan
$\binom{2(2l+1)}{N}$ determinan, dan \omterm{def:b3:quantum-model-systems:degenerate}{degenerasi} $(2L+1)(2S+1)$ suku-sukunya
berjumlah sama dengan bilangan itu.
\end{proposition}

\begin{proof}
Sebuah determinan ditentukan oleh himpunan \omterm{def:b3:many-electron-atoms:spin-orbital}{spin-orbital} yang terisi (urutannya
hanya mengubah tandanya): kita memilih $N$ dari $2(2l+1)$ \omterm{def:b3:many-electron-atoms:spin-orbital}{spin-orbital}. Suku-suku
itu adalah keadaan-keadaan yang sama yang dikelompokkan ulang, sehingga
jumlahnya harus sesuai.
\end{proof}

\begin{method}[Suku-suku sebuah konfigurasi]\label{met:b3:many-electron-atoms:terms}
\begin{enumerate}
  \item Daftarlah determinan-determinan yang diizinkan prinsip Pauli dan
        tabulasikan menurut $M_L$ dan $M_S$.
  \item Carilah $M_L$ terbesar; bersama $M_S$ terbesar yang menyertainya,
        nilai itu memulai sebuah suku dengan $L = M_L$, $S = M_S$.
  \item Coretlah satu determinan untuk setiap pasangan $(M_L, M_S)$ dengan
        $|M_L| \le L$, $|M_S| \le S$.
  \item Ulangi dengan sisanya sampai tabelnya kosong; periksalah jumlahnya.
\end{enumerate}
\end{method}

\begin{example}[Suku-suku $p^2$]\label{ex:b3:many-electron-atoms:p2}
Dua elektron dalam $p$ ($m_l = 1, 0, -1$) memberikan $\binom62 = 15$ determinan.
Tabelnya menurut $M_L$ dan $M_S$ adalah:
\begin{center}\small
\begin{tabular}{c|ccc}
\hline
 & $M_S = 1$ & $M_S = 0$ & $M_S = -1$\\ \hline
$M_L = 2$ & -- & $(1^+,1^-)$ & --\\
$M_L = 1$ & $(1^+,0^+)$ & $(1^+,0^-)$, $(1^-,0^+)$ & $(1^-,0^-)$\\
$M_L = 0$ & $(1^+,-1^+)$ & $(1^+,-1^-)$, $(1^-,-1^+)$, $(0^+,0^-)$ & $(1^-,-1^-)$\\
$M_L = -1$ & $(0^+,-1^+)$ & $(0^+,-1^-)$, $(0^-,-1^+)$ & $(0^-,-1^-)$\\
$M_L = -2$ & -- & $(-1^+,-1^-)$ & --\\ \hline
\end{tabular}
\end{center}
(dengan $m_l^\pm$ untuk $m_s = \pm\frac12$). $M_L = 2$ hanya muncul bersama
$M_S = 0$: sebuah suku ${}^1$D (5 keadaan). $M_L = 1$, yang terbesar di antara
sisanya, muncul bersama $M_S = 1$: sebuah suku ${}^3$P (9 keadaan). Tersisa
satu keadaan dengan $M_L = M_S = 0$: sebuah suku ${}^1$S. Memang
$5 + 9 + 1 = 15$.
\end{example}

\begin{example}[Karbon, terukur]\label{ex:b3:many-electron-atoms:carbon}
% ledger: lev:C.3P1, lev:C.3P2, lev:C.1D2, lev:C.1S0
Konfigurasi dasar $2p^2$ karbon memberikan tingkat-tingkat \termsym{3}{P}{0},
\termsym{3}{P}{1}, \termsym{3}{P}{2} pada 0, 16.4 dan \qty{43.4}{cm^{-1}},
kemudian \termsym{1}{D}{2} pada \qty{10193}{cm^{-1}} dan \termsym{1}{S}{0}
pada \qty{21648}{cm^{-1}}: suku ${}^3$P adalah yang terendah, seperti yang
diramalkan oleh aturan Hund.
\end{example}

\begin{omfigure}
\begin{tikzpicture}
\begin{axis}[name=main, width=0.62\linewidth, height=6.4cm, xmin=-0.3, xmax=6.6,
  ymin=-1500, ymax=23500, axis y line=left, axis x line=none,
  ylabel={energi / cm$^{-1}$}, unbounded coords=jump, clip=false,
  scaled y ticks=false, ytick={0,10000,20000},
  tick label style={font=\small}, label style={font=\small}]
\addplot[omstate] table[x=x, y=E] {figdata/out/bachelor-3-atomic-levels-carbon.dat};
\draw[densely dotted, gray] (axis cs:1,4858.5) -- (axis cs:2,29.6);
\draw[densely dotted, gray] (axis cs:1,4858.5) -- (axis cs:2,10192.7);
\draw[densely dotted, gray] (axis cs:1,4858.5) -- (axis cs:2,21648);
\draw[densely dotted, gray] (axis cs:3,10192.7) -- (axis cs:4,10192.7);
\draw[densely dotted, gray] (axis cs:3,21648) -- (axis cs:4,21648);
\node[font=\small, anchor=south] at (axis cs:0.5,4858.5) {$2p^2$};
\node[font=\small, anchor=south] at (axis cs:2.5,21648) {${}^1$S};
\node[font=\small, anchor=south] at (axis cs:2.5,10192.7) {${}^1$D};
\node[font=\small, anchor=south] at (axis cs:2.5,29.6) {${}^3$P};
\node[font=\small, anchor=west] at (axis cs:5.05,21648) {\termsym{1}{S}{0}};
\node[font=\small, anchor=west] at (axis cs:5.05,10192.7) {\termsym{1}{D}{2}};
\node[font=\scriptsize, anchor=north] at (axis cs:0.5,-600) {konfigurasi};
\node[font=\scriptsize, anchor=north] at (axis cs:2.5,-600) {suku};
\node[font=\scriptsize, anchor=north] at (axis cs:4.5,-600) {tingkat};
\draw[gray, ->] (axis cs:3.05,60) -- (axis cs:6.9,2500);
\end{axis}
\begin{axis}[at={(main.south east)}, anchor=south west, xshift=1.2cm, yshift=0.9cm,
  width=3.6cm, height=3.6cm, xmin=-0.1, xmax=1.1, ymin=-5, ymax=50,
  axis y line=right, axis x line=none, unbounded coords=jump, ytick={0,16.4,43.4},
  yticklabels={0,16.4,43.4}, tick label style={font=\scriptsize}, clip=false,
  title={\scriptsize ${}^3$P, diperbesar}, title style={yshift=-4pt}]
\addplot[omstate] table[x=x, y=E] {figdata/out/bachelor-3-atomic-levels-carbon-zoom.dat};
\node[font=\scriptsize, anchor=east] at (axis cs:0,0) {$J = 0$};
\node[font=\scriptsize, anchor=east] at (axis cs:0,16.4) {$J = 1$};
\node[font=\scriptsize, anchor=east] at (axis cs:0,43.4) {$J = 2$};
\end{axis}
\end{tikzpicture}

{\small Konfigurasi dasar karbon: satu konfigurasi, tiga suku (pada energi
terukurnya; suku ${}^3$P pada rata-rata berbobot tingkat-tingkatnya,
konfigurasi pada rata-rata kelima belas keadaannya) dan lima tingkat.
Pemisahan spin--orbit ${}^3$P (diperbesar, dalam \unit{cm^{-1}}) beberapa
ratus kali lebih kecil daripada jarak antarsuku.}
\end{omfigure}

\section{Aturan Hund dan kopling spin--orbit}

\begin{proposition}[Aturan Hund untuk suku]\label{prop:b3:many-electron-atoms:hund-terms}
Untuk konfigurasi dasar sebuah atom atau ion, suku terendah adalah suku
dengan (1) $S$ terbesar, lalu (2) untuk $S$ itu, $L$ terbesar. Tingkat
terendahnya mempunyai (3) $J = |L - S|$ jika subkulitnya kurang dari setengah
penuh, dan $J = L +
S$ jika lebih dari setengah penuh.
\end{proposition}
\begin{proof}[Kedudukan aturan ini]
Aturan-aturan ini ditetapkan secara eksperimen, untuk konfigurasi dasar
saja; aturan-aturan itu bukan teorema. Aturan 1 adalah stabilisasi
pertukaran dari subbab sebelumnya; aturan 2 menyatakan bahwa elektron-elektron
yang berputar searah lebih jarang bertemu; aturan 3 mengikuti tanda \omterm{def:b3:many-electron-atoms:spin-orbit}{kopling
spin--orbit}, di bawah ini.
\end{proof}

\begin{method}[Suku dasar dari diagram kotak]\label{met:b3:many-electron-atoms:ground-term}
\begin{enumerate}
  \item Isilah kotak-kotak subkulit terbuka mulai dari $m_l = +l$ ke bawah,
        satu elektron per kotak dengan spin sejajar, lalu pasangkan lagi
        mulai dari $+l$.
  \item $S = \frac12 \times$(banyaknya elektron tak berpasangan); $L = |\sum m_l|$.
  \item $J = |L - S|$ untuk kurang dari setengah penuh, $L + S$ untuk lebih
        dari setengah penuh; untuk tepat setengah penuh, $L = 0$ dan $J = S$.
\end{enumerate}
\end{method}

\begin{example}[Suku-suku dasar]\label{ex:b3:many-electron-atoms:ground-terms}
% ledger: lev:Ti2.3P0, lev:Ni2.3P0, lev:Cr3.4P1_2
\ce{Ti^2+}, $d^2$: kotak $m_l = 2, 1$ terisi, $S = 1$, $L = 3$, $J = 2$:
\termsym{3}{F}{2}. \ce{Cr^3+}, $d^3$: $S = \frac32$, $L = 2 + 1 + 0 = 3$,
\termsym{4}{F}{3/2}. \ce{Fe^2+}, $d^6$: lima ke atas, satu ke bawah di
$m_l = 2$: $S = 2$, $L = 2$, lebih dari setengah penuh, \termsym{5}{D}{4}.
\ce{Ni^2+}, $d^8$: $S = 1$, $L = 3$, \termsym{3}{F}{4}. Masing-masing sesuai
dengan tingkat dasar yang diukur spektroskopi atom.
\end{example}

\begin{definition}[Kopling spin--orbit, struktur halus]\label{def:b3:many-electron-atoms:spin-orbit}
Yang disebut \emph{kopling spin--orbit}\index{kopling spin--orbit} adalah
interaksi antara momen magnetik spin sebuah elektron dan momen magnetik gerak
orbitalnya, yang untuk sebuah suku ditulis $A\,\hat{\mathbf L}\cdot\hat{\mathbf S}/\hbar^2$.
Kopling ini memecah sebuah suku menjadi tingkat-tingkatnya dengan $J$ yang
berbeda: itulah \emph{struktur halus}\index{struktur halus} suku tersebut.
\end{definition}

\begin{proposition}[Aturan interval Landé]\label{prop:b3:many-electron-atoms:lande}
Di dalam sebuah suku, $E(J) = E_0 + \frac{A}{2}[J(J+1) - L(L+1) - S(S+1)]$,
sehingga $E(J) - E(J - 1) = AJ$. $A > 0$ (normal) untuk subkulit yang kurang
dari setengah penuh, $A < 0$ (terbalik) untuk subkulit yang lebih dari
setengah penuh.
\end{proposition}

\begin{proof}
$\hat{\mathbf J} = \hat{\mathbf L} + \hat{\mathbf S}$ memberikan $\hat{\mathbf J}^2 =
\hat{\mathbf L}^2 + \hat{\mathbf S}^2 + 2\hat{\mathbf L}\cdot\hat{\mathbf S}$, sehingga
$\hat{\mathbf L}\cdot\hat{\mathbf S} = \frac12(\hat{\mathbf J}^2 - \hat{\mathbf L}^2 -
\hat{\mathbf S}^2)$, yang \omterm{def:b3:quantum-model-systems:operator}{nilai eigennya} dalam keadaan dengan $J$, $L$, $S$
tertentu adalah $\frac{\hbar^2}{2}[J(J+1) - L(L+1) - S(S+1)]$. Selisih antara
$J$ dan $J - 1$ adalah $\frac A2[J(J+1) - (J-1)J] = AJ$. Tanda $A$ diterima
tanpa bukti: subkulit yang lebih dari setengah penuh berperilaku seperti
sejumlah lubang positif yang bersesuaian.
\end{proof}

\begin{example}[Menguji aturan interval]\label{ex:b3:many-electron-atoms:lande-test}
% ledger: lev:C.3P1, lev:C.3P2, lev:O.3P1, lev:O.3P0
Untuk ${}^3$P, aturan itu meramalkan interval dengan rasio $J = 2 : 1$, yaitu 2.
Karbon: $(43.4 - 16.4)/16.4 = 1.65$. Oksigen ($2p^4$, terbalik, \termsym{3}{P}{2}
terendah, lalu $J = 1$ pada \qty{158.3}{cm^{-1}} dan $J = 0$ pada
\qty{227.0}{cm^{-1}}): $158.3/68.7 = 2.30$. Aturan itu berlaku dengan selisih
kurang dari 20\,\%: \omterm{def:b3:many-electron-atoms:russell-saunders}{kopling Russell--Saunders} adalah deskripsi yang baik untuk
atom ringan, bukan deskripsi yang eksak.
\end{example}

\begin{omfigure}
\begin{tikzpicture}[font=\small]
  \draw[omstate] (0,0) -- (1.6,0) node[midway, below] {$3s\ {}^2$S$_{1/2}$};
  \draw[omstate] (0,3.0) -- (1.6,3.0);
  \node[anchor=east] at (-0.1,3.0) {$3p$};
  \draw[omstate] (2.8,2.75) -- (4.4,2.75) node[right] {${}^2$P$_{1/2}$};
  \draw[omstate] (2.8,3.25) -- (4.4,3.25) node[right] {${}^2$P$_{3/2}$};
  \draw[densely dotted, gray] (1.6,3.0) -- (2.8,2.75) (1.6,3.0) -- (2.8,3.25);
  \draw[omfluo] (3.3,2.73) -- (3.3,0.02) node[pos=0.55, left] {D$_1$};
  \draw[omfluo] (3.9,3.23) -- (3.9,0.02) node[pos=0.55, right] {D$_2$};
  \draw[omstate] (2.8,0) -- (4.4,0);
  \draw[densely dotted, gray] (1.6,0) -- (2.8,0);
  \draw[decorate, decoration={brace, amplitude=3pt}] (5.5,3.27) -- (5.5,2.73)
     node[midway, right=3pt, align=left] {\footnotesize \qty{17.2}{cm^{-1}}\\[-1pt]\footnotesize (tidak berskala)};
\end{tikzpicture}

{\small Garis-garis D natrium. Tingkat $3p$ terpecah oleh \omterm{def:b3:many-electron-atoms:spin-orbit}{kopling spin--orbit}
menjadi ${}^2$P$_{1/2}$ dan ${}^2$P$_{3/2}$; emisi ke tingkat dasar $3s$
memberikan dua garis kuning, D$_1$ pada \qty{589.76}{nm} dan D$_2$ pada
\qty{589.16}{nm} (panjang gelombang dalam vakum).}
\end{omfigure}

\begin{example}[Doblet natrium]\label{ex:b3:many-electron-atoms:sodium}
% ledger: lev:Na.3p1_2, lev:Na.3p3_2
Kedua tingkat $3p$ natrium berada pada 16956.17 dan \qty{16973.37}{cm^{-1}}:
garis-garis D jatuh pada $10^7/16956.17 = \qty{589.76}{nm}$ dan
\qty{589.16}{nm} dalam vakum, terpisah sejauh \qty{17.20}{cm^{-1}}. Untuk
${}^2$P, $E(\frac32) - E(\frac12) =
\frac32A$, sehingga $A = \qty{11.47}{cm^{-1}}$.
\end{example}

\begin{proposition}[Aturan seleksi untuk atom]\label{prop:b3:many-electron-atoms:selection-atoms}
Dalam \omterm{def:b3:many-electron-atoms:russell-saunders}{kopling Russell--Saunders}, sebuah transisi dipol listrik mensyaratkan
$\Delta S = 0$, $\Delta L = 0, \pm1$, $\Delta J = 0, \pm1$ (tetapi bukan $J = 0 \to J
= 0$), dan perubahan paritas: untuk loncatan satu elektron, $\Delta l = \pm1$.
\end{proposition}
\begin{proof}[Kedudukan aturan ini]
$\Delta S = 0$ dibuktikan dalam \cref{ch:b3:electronic-spectroscopy}: \omterm{def:b3:quantum-model-systems:operator}{operator}
dipol tidak bekerja pada spin. Aturan-aturan lainnya mengikuti sifat vektor
\omterm{def:b3:quantum-model-systems:operator}{operator} dipol dan di sini diterima tanpa bukti.
\end{proof}

\begin{omfigure}
\begin{tikzpicture}
\begin{axis}[width=0.62\linewidth, height=6.2cm, xmin=-0.5, xmax=3.5,
  ymin=158, ymax=173, axis y line=left, axis x line=none,
  ylabel={energi / $10^3$ cm$^{-1}$}, unbounded coords=jump, clip=false,
  tick label style={font=\small}, label style={font=\small}]
\addplot[omstate] table[x=x, y=E] {figdata/out/bachelor-3-atomic-levels-helium.dat};
\node[font=\small, anchor=west] at (axis cs:1.05,166.277) {$1s2s\ {}^1$S};
\node[font=\small, anchor=west] at (axis cs:1.05,171.135) {$1s2p\ {}^1$P};
\node[font=\small, anchor=west] at (axis cs:3.05,159.856) {$1s2s\ {}^3$S};
\node[font=\small, anchor=west] at (axis cs:3.05,169.087) {$1s2p\ {}^3$P};
\node[font=\small, anchor=south] at (axis cs:0.5,172.3) {singlet};
\node[font=\small, anchor=south] at (axis cs:2.5,172.3) {triplet};
\draw[omfluo] (axis cs:0.5,171.0) -- (axis cs:0.5,166.42) node[pos=0.5, left, font=\scriptsize] {\qty{2059}{nm}};
\draw[omfluo] (axis cs:2.5,168.95) -- (axis cs:2.5,160.0) node[pos=0.5, left, font=\scriptsize] {\qty{1083}{nm}};
\draw[densely dotted, gray] (axis cs:-0.3,159.856) -- (axis cs:2,159.856);
\draw[<->, omThm] (axis cs:-0.2,159.856) -- (axis cs:-0.2,166.277) node[pos=0.5, right, font=\scriptsize] {$2K$};
\end{axis}
\end{tikzpicture}

{\small Dua ``helium''. Tingkat-tingkat singlet dan triplet dari konfigurasi
$1s2s$ dan $1s2p$ pada energi terukurnya di atas keadaan dasar (suku ${}^3$P
pada rata-rata tingkat-tingkatnya). Garis-garis hanya menghubungkan
keadaan-keadaan dari keluarga yang sama; setiap triplet berada $2K$ di bawah
singletnya.}
\end{omfigure}

\begin{history}[Dua helium dan prinsip larangan]
Helium dinamai menurut Matahari pada tahun 1868 dan ditemukan di Bumi oleh
William Ramsay pada tahun 1895. Kedua deret garisnya membingungkan para
spektroskopis selama tiga puluh tahun. Pada tahun 1925 Wolfgang Pauli
mengusulkan bahwa tidak ada dua elektron yang memiliki bilangan kuantum yang
sama, pada tahun yang sama ketika spin elektron diusulkan oleh George
Uhlenbeck dan Samuel Goudsmit; pada tahun 1926 Werner Heisenberg menunjukkan
bahwa simetri fungsi gelombang saja sudah memecah helium menjadi keluarga
singlet dan keluarga triplet.
\end{history}

\begin{inthelab}[Memisahkan doblet natrium]
Sebuah lampu natrium diletakkan di depan celah spektrometer kisi. Dengan
kisi 600 garis per milimeter, sudut orde pertama kedua garis itu berbeda
sekitar \num{0.02}\textdegree: spektrometer yang baik memisahkannya, prisma
praktikum mahasiswa tidak. Lampu itu menjadi panas; lampu hanya dipegang
setelah dingin.
\end{inthelab}

\section{Latihan}

\begin{exercise}[$\star$]\label{exo:b3:many-electron-atoms:1}
Tuliskan \omterm{def:b3:many-electron-atoms:slater}{determinan Slater} untuk keadaan dasar litium, $1s^22s$, dengan
elektron $2s$ berspin $\alpha$. Mengapa tidak ada determinan untuk
$1s^3$?
\end{exercise}

\begin{exercise}[$\star$]\label{exo:b3:many-electron-atoms:2}
Hitunglah banyaknya determinan konfigurasi $p^2$, $p^3$ dan $d^2$. Suku-suku
$p^3$ adalah ${}^4$S, ${}^2$D, ${}^2$P dan suku-suku $d^2$ adalah ${}^3$F, ${}^3$P, ${}^1$G,
${}^1$D, ${}^1$S: periksalah kedua penghitungan itu.
\end{exercise}

\begin{exercise}[$\star$]\label{exo:b3:many-electron-atoms:3}
Berikan suku dan tingkat dasar N, O, F, \ce{Ti^2+} dan \ce{Fe^3+} ($d^5$).
\end{exercise}

\begin{exercise}[$\star$]\label{exo:b3:many-electron-atoms:4}
Daftarlah tingkat-tingkat suku ${}^2$D dan ${}^3$P beserta \omterm{def:b3:quantum-model-systems:degenerate}{degenerasinya} $2J+1$,
dan periksalah bahwa jumlah \omterm{def:b3:quantum-model-systems:degenerate}{degenerasinya} sama dengan $(2L+1)(2S+1)$.
\end{exercise}

\begin{exercise}[$\star\star$]\label{exo:b3:many-electron-atoms:5}
Tingkat-tingkat ${}^3$P karbon terletak pada 0, 16.4 dan \qty{43.4}{cm^{-1}}.
Hitunglah rasio interval-intervalnya dan tetapan spin--orbit $A$ dari setiap
interval. Apa arti perbedaan antara kedua nilai $A$ itu?
\end{exercise}

\begin{exercise}[$\star\star$]\label{exo:b3:many-electron-atoms:6}
Tingkat-tingkat $3p$ natrium terletak pada 16956.17 dan \qty{16973.37}{cm^{-1}}.
Hitunglah panjang gelombang garis-garis D dalam vakum, jaraknya dalam
\unit{cm^{-1}} dan dalam meV, serta $A$ untuk suku $3p$.
\end{exercise}

\begin{exercise}[$\star\star$]\label{exo:b3:many-electron-atoms:7}
Manakah di antara transisi natrium berikut yang diizinkan dalam emisi:
$3p \to 3s$, $3d \to 3s$, $3d \to 3p$, $4s \to 3s$, $4s \to 3p$? Berikan
alasan untuk masing-masing.
\end{exercise}

\begin{exercise}[$\star\star$]\label{exo:b3:many-electron-atoms:8}
% ledger: lev:He.2p3P2, lev:He.2p3P1, lev:He.2p3P0, lev:He.2p1P1
Tingkat-tingkat $1s2p$ helium adalah \termsym{3}{P}{2}, \termsym{3}{P}{1},
\termsym{3}{P}{0} pada 169086.76, 169086.84, \qty{169087.83}{cm^{-1}} dan
\termsym{1}{P}{1} pada \qty{171134.89}{cm^{-1}}. Hitunglah energi rata-rata
suku ${}^3$P, lalu $K(1s,2p)$ dalam eV. Bandingkan dengan $K(1s,2s) = \qty{0.398}{eV}$ dan jelaskan perbedaannya.
\end{exercise}

\begin{exercise}[$\star\star$]\label{exo:b3:many-electron-atoms:9}
Carilah suku-suku $d^2$ dengan metode bab ini, mulai dari $M_L$ terbesar
(seluruh tabel tidak perlu ditulis, tetapi berikan alasan untuk setiap suku).
\end{exercise}

\begin{exercise}[$\star\star\star$]\label{exo:b3:many-electron-atoms:10}
Tunjukkan bahwa fungsi ruang triplet $\frac{1}{\sqrt2}[1s(1)2s(2) -
2s(1)1s(2)]$ lenyap ketika $\mathbf r_1 = \mathbf r_2$ sedangkan fungsi singlet
tidak. Kaitkan hal ini dengan tanda $E_+ - E_-$.
\end{exercise}

\begin{exercise}[$\star\star\star$]\label{exo:b3:many-electron-atoms:11}
% ledger: lev:Ni2.3F3, lev:Ni2.3F2
Untuk \ce{Ni^2+} ($d^8$), tingkat-tingkat ${}^3$F terletak pada 0 ($J = 4$),
1360.7 ($J = 3$) dan \qty{2269.6}{cm^{-1}} ($J = 2$). Jelaskan urutannya,
hitunglah rasio interval-intervalnya, dan bandingkan dengan ramalan Landé.
Berikan $A$ dari interval yang lebih besar.
\end{exercise}

\begin{exercise}[$\star\star\star$]\label{exo:b3:many-electron-atoms:12}
Tunjukkan bahwa \omterm{def:b3:quantum-model-systems:hermitian}{nilai harapan} $\hat{\mathbf L}\cdot\hat{\mathbf S}$, yang
dijumlahkan atas semua $(2L+1)(2S+1)$ keadaan sebuah suku, bernilai nol,
sehingga \omterm{def:b3:many-electron-atoms:spin-orbit}{kopling spin--orbit} tidak mengubah rata-rata berbobot
tingkat-tingkatnya. Periksalah hal ini pada ${}^3$P karbon.
\end{exercise}

\section{Soal: Dua Helium}

\begin{problem}[{Soal akhir pekan --- keluarga singlet dan triplet helium:
determinan, suku-suku ion logam transisi, aturan seleksi yang memisahkan
kedua keluarga, dan integral pertukaran yang terukur}]
\label{pb:b3:many-electron-atoms:1}
% ledger: lev:He.2s3S1, lev:He.2s1S0, lev:He.2p3P2, lev:He.2p3P1, lev:He.2p3P0, lev:He.2p1P1, lev:Ti2.3F3, lev:Ti2.3F4, lev:Ti2.3P0, lev:Ti2.3P1, lev:Ti2.3P2
Data (NIST, dalam \unit{cm^{-1}} di atas keadaan dasar setiap spesies). Helium:
$1s2s$ \termsym{3}{S}{1} 159855.97, \termsym{1}{S}{0} 166277.44; $1s2p$
\termsym{3}{P}{2,1,0} 169086.76, 169086.84, 169087.83, \termsym{1}{P}{1}
171134.89. \ce{Ti^2+}: \termsym{3}{F}{2,3,4} 0, 184.9, 420.4; \termsym{3}{P}{0,1,2}
10538.4, 10603.6, 10721.2. $\qty{1}{eV} = \qty{8065.54}{cm^{-1}}$.

\medskip\noindent\textbf{Bagian I --- \omterm{def:b3:many-electron-atoms:spin-orbital}{Spin-orbital} dan determinan.}
\begin{enumerate}[beginpenalty=10000]
  \item Daftarlah keempat \omterm{def:b3:many-electron-atoms:spin-orbital}{spin-orbital} yang tersedia bagi konfigurasi $1s2s$.
  \item Tuliskan keempat determinan dengan satu elektron dalam $1s$ dan satu
        dalam $2s$.
  \item Tunjukkan bahwa $|1s\alpha\,2s\alpha|$ adalah hasil kali sebuah fungsi
        ruang antisimetris dan fungsi spin $\alpha(1)\alpha(2)$.
  \item Gabungkan $|1s\alpha\,2s\beta|$ dan $|1s\beta\,2s\alpha|$ untuk
        memperoleh komponen $M_S = 0$ dari fungsi ruang yang sama, dan sebuah
        keadaan keempat.
  \item Fungsi ruang mana, simetris atau antisimetris, yang menyertai
        $S = 0$? Dan $S = 1$?
  \item Mengapa keadaan-keadaan ini disebut singlet dan triplet?
\end{enumerate}

\medskip\noindent\textbf{Bagian II --- Suku-suku \ce{Ti^2+}.}
\begin{enumerate}[resume, beginpenalty=10000]
  \item Berikan konfigurasi \ce{Ti^2+} dan banyaknya determinannya.
  \item Carilah suku dan tingkat dasarnya dengan aturan Hund; periksalah
        dengan data.
  \item Suku-suku lainnya adalah ${}^3$P, ${}^1$G, ${}^1$D dan ${}^1$S: periksalah
        jumlah keadaannya.
  \item Apakah \omterm{def:b3:many-electron-atoms:spin-orbit}{struktur halus} ${}^3$F normal atau terbalik? Ujilah aturan
        interval Landé.
  \item Hitunglah energi rata-rata berbobot suku ${}^3$F dan ${}^3$P.
  \item Jarak antara keduanya adalah $15B$, dengan $B$ sebuah parameter Racah
        ion itu (\cref{ch:b3:complex-spectra-magnetism}). Hitunglah $B$.
\end{enumerate}

\medskip\noindent\textbf{Bagian III --- Mengapa ada dua keluarga.}
\begin{enumerate}[resume, beginpenalty=10000]
  \item Aturan seleksi mana yang melarang transisi antara singlet dan
        triplet?
  \item Triplet terendah, $1s2s$ \termsym{3}{S}{1}, tidak dapat memancar ke
        keadaan dasar $1s^2$ \termsym{1}{S}{0}. Berikan dua alasan. Disebut
        apakah keadaan semacam itu?
  \item Hitunglah panjang gelombang garis $1s2p\ {}^3$P $\to 1s2s\ {}^3$S
        (gunakan rata-rata tingkat-tingkat ${}^3$P).
  \item Hitunglah panjang gelombang $1s2p\ {}^1$P $\to 1s2s\ {}^1$S.
  \item Hitunglah panjang gelombang $1s2p\ {}^1$P $\to 1s^2\ {}^1$S. Di daerah
        spektrum mana garis itu berada?
  \item Jelaskan mengapa para spektroskopis abad kesembilan belas, ketika
        melihat garis-garis ini, percaya bahwa ada dua gas yang berbeda.
\end{enumerate}

\medskip\noindent\textbf{Bagian IV --- \omterm{def:b3:many-electron-atoms:exchange}{Integral pertukaran}, terukur.}
\begin{enumerate}[resume, beginpenalty=10000]
  \item Tuliskan energi orde pertama singlet dan triplet $1s2s$ dalam
        $E_{1s}$, $E_{2s}$, $J$ dan $K$.
  \item Manakah yang lebih rendah, dan mengapa, ditinjau dari posisi
        elektron-elektronnya?
  \item Hitunglah selisih singlet--triplet $1s2s$ dalam \unit{cm^{-1}} dan eV.
  \item Hitunglah $K(1s,2p)$ dalam eV dari tingkat-tingkat $1s2p$.
  \item Mengapa $K(1s,2s)$ lebih besar daripada $K(1s,2p)$?
  \item Apakah aturan Hund yang pertama didukung oleh keadaan-keadaan
        tereksitasi helium?
  \item Nyatakan hasilnya: \omterm{def:b3:many-electron-atoms:exchange}{integral pertukaran} $K(1s,2s)$ helium, dalam eV.
\end{enumerate}
\end{problem}
